% ---
% titre: Priorité des opération
% theme: NO
% chapitre: Nombres relatifs
% sous_chapitre: Priorité des opérations
% niveau: [10e]
% type: EVAL
% tags: [c-priorite-operations,c-puissance,c-racine,c-nombres-relatifs,p-question-TS]
% difficulte: 2
% corrige: true
% ---

\question[19]
Calcule le résultat des opérations suivantes \textbf{en écrivant les étapes intermédiaires.}

\vspace{30pt}

\begin{parts}
	\part
		$24 + 6 \cdot (13-7)=$
		
	\vspace{10pt}{\footnotesize\raggedleft Espace pour ta démarche et tes calculs\par}
	\begin{solutionorgrid}[100pt]
	
		\vspace{10pt}
		$24 + 6 \cdot 6 \textbf{\,(0.5pt)} = 24 + 36 \textbf{\,(0.5pt)} = 60\textbf{\,(0.5pt)}$
	\end{solutionorgrid}	
	
	\vfill	
	\part
		$79 - \sqrt{25} \cdot (15-6)=$
		
	\vspace{10pt}{\footnotesize\raggedleft Espace pour ta démarche et tes calculs\par}
	\begin{solutionorgrid}[100pt]
	
		\vspace{10pt}
		$79 -5 \textbf{\,(0.5pt)} \cdot (15-6) =  79 -5 \cdot 9\textbf{\,(0.5pt)} = 79 - 45 \textbf{\,(0.5pt)}  = 34 \textbf{\,(0.5pt)} $
	\end{solutionorgrid}
	
	\vfill	
	\part
		$4^3 + 37 - (30:10)^2=$
		
	\vspace{10pt}{\footnotesize\raggedleft Espace pour ta démarche et tes calculs\par}
	\begin{solutionorgrid}[100pt]
	
		\vspace{10pt}
		$64 \textbf{\,(0.5pt)} + 37 - 3^2 \textbf{\,(0.5pt)} = 64 + 37 - 9\textbf{\,(0.5pt)}  = 92$
	\end{solutionorgrid}	
	
	\vfill	
	\part
		$7,6 + 3,4 + 18 \cdot 0,1=$
		
	\vspace{10pt}{\footnotesize\raggedleft Espace pour ta démarche et tes calculs\par}
	\begin{solutionorgrid}[100pt]
	
		\vspace{10pt}
		$7,6 + 3,4 + 1,8 \textbf{\,(0.5pt)}  = 12,8 \textbf{\,(0.5pt)} $
	\end{solutionorgrid}
		
	\newpage
	\part
		$(-5,6) - (-2,2) + (+12,6)$
		
	\vspace{10pt}{\footnotesize\raggedleft Espace pour ta démarche et tes calculs\par}
	\begin{solutionorgrid}[100pt]
	
		\vspace{10pt}
		$-5,6 + 2,2 + 12,6 \textbf{\,(0.5pt)}  = 14,8 - 5,6 \textbf{\,(0.5pt)}  = 9,2 \textbf{\,(0.5pt)} $
	\end{solutionorgrid}	
		
	\vfill
	\part
		$(-6) \cdot (-7) : \sqrt{16}=$
		
	\vspace{10pt}{\footnotesize\raggedleft Espace pour ta démarche et tes calculs\par}
	\begin{solutionorgrid}[100pt]
	
		\vspace{10pt}
		$(-6) \cdot (-7) : 4 \textbf{\,(0.5pt)} = 42 \textbf{\,(0.5pt)} :4 = 10,5 \textbf{\,(0.5pt)} $
	\end{solutionorgrid}	
	
	\vfill	
	\part
		$(-15) \cdot (-12) \cdot (-5)=$
		
	\vspace{10pt}{\footnotesize\raggedleft Espace pour ta démarche et tes calculs\par}
	\begin{solutionorgrid}[100pt]
	
		\vspace{10pt}
		$- 900\textbf{\,(1pt)}$
	\end{solutionorgrid}			

	\vfill
	\part
		$(-100):(-20)\cdot(-5)^2=$
		
	\vspace{10pt}{\footnotesize\raggedleft Espace pour ta démarche et tes calculs\par}
	\begin{solutionorgrid}[100pt]
	
		\vspace{10pt}
		$(-100):(-20)\cdot25\textbf{\,(0.5pt)} =5 \textbf{\,(0.5pt)}  \cdot 25 = 125 \textbf{\,(0.5pt)} $
	\end{solutionorgrid}	
	
	\vfill
	\part
		$88 + 12 \cdot 11 - 1=$
		
	\vspace{10pt}{\footnotesize\raggedleft Espace pour ta démarche et tes calculs\par}
	\begin{solutionorgrid}[100pt]
	
		\vspace{10pt}
		$88 + 132 \textbf{\,(0.5pt)}  - 1 = 219 \textbf{\,(0.5pt)} $
	\end{solutionorgrid}	
	
	\vfill
	\part
		$(+9) \cdot (-5)^2 - (-25)=$
		
	\vspace{10pt}{\footnotesize\raggedleft Espace pour ta démarche et tes calculs\par}
	\begin{solutionorgrid}[100pt]
	
		\vspace{10pt}
		$9 \cdot 25 \textbf{\,(0.5pt)} - (-25) = 225 \textbf{\,(0.5pt)} - (-25)= 225 +25 \textbf{\,(0.5pt)} = 250 \textbf{\,(0.5pt)}$
	\end{solutionorgrid}	
	
	\vfill
	\part
		$(1 + 2 \cdot 3)^2 \cdot \sqrt{\left[(-11) \cdot (-11) - 3 \cdot 7\right]}-1=$
		
	\vspace{10pt}{\footnotesize\raggedleft Espace pour ta démarche et tes calculs\par}
	\begin{solutionorgrid}[100pt]
	
		\vspace{10pt}
		$(1 + 6\textbf{\,(0.5pt)})^2 \cdot \sqrt{\left(121\textbf{\,(0.5pt)} -21 \textbf{\,(0.5pt)}\right)}-1 = 7^2\textbf{\,(0.5pt)} \cdot \sqrt{100 \textbf{\,(0.5pt)}}-1  =  49\textbf{\,(0.5pt)} \cdot 10 \textbf{\,(0.5pt)}-1  = 490 \textbf{\,(0.5pt)} -1 = 489 \textbf{\,(0.5pt)} $
	\end{solutionorgrid}		
	\vfill
\end{parts}
